\chapter{Dám Nói Không}
\markboth{Dám Nói Không}{The Courage to Say No}

\section{Biết Nói Không Là Một Dấu Hiệu Trưởng Thành}
\begin{truyen}{Biết Nói Không Là Một Dấu Hiệu Trưởng Thành}{Knowing How to Say No Is a Sign of Maturity}
\chuhoa{N}{hiều học sinh nghĩ từ chối là hỗn, là ích kỷ, là làm mất lòng.}
\textit{(Many students think refusal is rude, selfish, or hurtful.)}

Vì vậy các em gật đầu với việc chép bài hộ, đi chơi khi không muốn, tham gia trò ngu khi thấy sai, và nhận thêm gánh khi đã quá sức.
\textit{(So they say yes to doing others’ homework, going out when they do not want to, joining stupid acts they know are wrong, and taking on burdens when already overloaded.)}

Một chữ không nói ra đúng lúc thường dẫn tới hàng chục hậu quả phải trả sau đó.
\textit{(One no withheld at the right time often leads to dozens of consequences paid later.)}

Từ chối không làm em thành người xấu.
\textit{(Refusing does not make you a bad person.)}

Từ chối đúng lúc giúp em giữ thời gian, nhân cách, và sự an toàn của mình.
\textit{(Refusing at the right time protects your time, your character, and your safety.)}

\end{truyen}

\begin{baihoc}
Giáo dục tốt không chỉ dạy học sinh vâng lời. Nó còn phải dạy các em từ chối điều sai và điều quá sức.
\textit{(Good education does not only teach students to obey. It must also teach them to refuse what is wrong and what is too much.)}
\end{baihoc}

\begin{ghinhoanh}
\textbf{Short English to remember:}

Good education does not only teach students to obey.

It must also teach them to refuse what is wrong and what is too much.

\end{ghinhoanh}

\begin{tuhocnhanh}
\textbf{Cau hoi tu hoc / Self-study questions}

1. Van de chinh la gi? \textit{What is the main problem?}

2. Cach xu ly tot hon la gi? \textit{What is a better action?}

3. Mot cau tieng Anh em muon nho la cau nao? \textit{Which English sentence do you want to remember?}
\end{tuhocnhanh}
\ngancach

\section{Ai Cũng Muốn Em Dễ Sai Khi Em Không Có Ranh Giới}
\begin{truyen}{Ai Cũng Muốn Em Dễ Sai Khi Em Không Có Ranh Giới}{Everyone Can Push You Around If You Have No Boundaries}
\chuhoa{M}{ột học sinh không có ranh giới rõ thường trở thành người bị nhờ nhiều nhất, lợi dụng nhiều nhất, và mệt mỏi nhiều nhất.}
\textit{(A student without clear boundaries often becomes the one most asked, most used, and most exhausted.)}

Không phải vì em xấu số.
\textit{(This is not bad luck.)}

Mà vì những người khác sớm nhận ra em khó từ chối.
\textit{(It is because others quickly notice that you are hard to refuse.)}

Ranh giới không cần hung hăng.
\textit{(Boundaries do not need aggression.)}

Nó chỉ cần rõ ràng: em không cho phép gì, em sẵn sàng gì, và em dừng ở đâu.
\textit{(They only need clarity: what you do not allow, what you are willing to give, and where you stop.)}

Người quen lợi dụng em sẽ khó chịu khi em bắt đầu có ranh giới.
\textit{(People who benefit from using you will dislike your boundaries.)}

Điều đó không chứng minh em sai.
\textit{(That does not prove you are wrong.)}

Nó chỉ chứng minh trước đây họ được lợi từ việc em quá mềm.
\textit{(It only proves they previously benefited from your softness.)}

\end{truyen}

\begin{baihoc}
Ranh giới là kỹ năng tự bảo vệ, không phải thái độ thù địch. Học sinh càng tử tế càng cần được học điều này.
\textit{(Boundaries are a self-protection skill, not hostility. The kinder the student, the more necessary this lesson becomes.)}
\end{baihoc}

\begin{ghinhoanh}
\textbf{Short English to remember:}

Boundaries are a self-protection skill, not hostility.

The kinder the student, the more necessary this lesson becomes.

\end{ghinhoanh}

\begin{tuhocnhanh}
\textbf{Cau hoi tu hoc / Self-study questions}

1. Van de chinh la gi? \textit{What is the main problem?}

2. Cach xu ly tot hon la gi? \textit{What is a better action?}

3. Mot cau tieng Anh em muon nho la cau nao? \textit{Which English sentence do you want to remember?}
\end{tuhocnhanh}
\ngancach

\section{Từ Chối Bạn Xấu Khó Hơn Từ Chối Người Lạ}
\begin{truyen}{Từ Chối Bạn Xấu Khó Hơn Từ Chối Người Lạ}{It Is Harder to Refuse Bad Friends Than Strangers}
\chuhoa{N}{hiều học sinh không sa vào rắc rối vì người lạ ép.}
\textit{(Many students do not get into trouble because strangers pressure them.)}

Các em sa vào vì bạn thân rủ.
\textit{(They get into trouble because close friends invite them.)}

Một điếu thuốc đầu tiên, một lần quay cóp, một trò làm nhục người khác, một cú trốn học.
\textit{(A first cigarette, one act of cheating, humiliating someone else, skipping school.)}

Những bước lệch lớn rất hay bắt đầu bằng câu: ``Có gì đâu, đi một lần thôi.''
\textit{(Major deviations often begin with the line: ``It is nothing, just once.'')}

Từ chối người mình quý rất khó vì em sợ mất nhóm, sợ bị chê là nhát, sợ trở thành kẻ khác biệt.
\textit{(Refusing people you like is hard because you fear losing the group, being mocked as weak, or becoming the odd one out.)}

Nhưng nhiều cuộc đời hỏng không phải vì một quyết định lớn.
\textit{(But many lives are not ruined by one huge decision.)}

Nó hỏng vì quá nhiều lần không dám nói không với cái sai nhỏ.
\textit{(They are ruined by too many small moments of failing to say no to what is wrong.)}

\end{truyen}

\begin{baihoc}
Hãy tập cho học sinh những câu từ chối cụ thể. Bản lĩnh không tự xuất hiện trong khoảnh khắc áp lực nếu trước đó chưa từng luyện.
\textit{(Train students with concrete refusal sentences. Strength does not appear magically under pressure if it has never been practiced before.)}
\end{baihoc}

\begin{ghinhoanh}
\textbf{Short English to remember:}

Train students with concrete refusal sentences.

Strength does not appear magically under pressure if it has never been practiced before.

\end{ghinhoanh}

\begin{tuhocnhanh}
\textbf{Cau hoi tu hoc / Self-study questions}

1. Van de chinh la gi? \textit{What is the main problem?}

2. Cach xu ly tot hon la gi? \textit{What is a better action?}

3. Mot cau tieng Anh em muon nho la cau nao? \textit{Which English sentence do you want to remember?}
\end{tuhocnhanh}
\ngancach

\section{Nói Không Với Chính Mình Còn Khó Hơn}
\begin{truyen}{Nói Không Với Chính Mình Còn Khó Hơn}{Saying No to Yourself Is Even Harder}
\chuhoa{K}{hông phải mọi cám dỗ đều đến từ người khác.}
\textit{(Not every temptation comes from other people.)}

Rất nhiều lần, chính mình dụ mình.
\textit{(Very often, we tempt ourselves.)}

Thêm mười phút điện thoại, thêm một lần trì hoãn, thêm một lời cay nghiệt cho đã miệng, thêm một lý do để không làm điều nên làm.
\textit{(Ten more minutes on the phone, one more delay, one more cruel remark for satisfaction, one more excuse not to do what should be done.)}

Tự chủ bắt đầu từ những khoảnh khắc rất nhỏ đó.
\textit{(Self-control begins in those tiny moments.)}

Người nói không được với chính mình thì sớm muộn cũng khó đứng vững trước sức ép bên ngoài.
\textit{(A person who cannot refuse themselves will eventually struggle to stand against outside pressure too.)}

Trưởng thành không chỉ là tự do làm điều mình muốn.
\textit{(Maturity is not only the freedom to do what you want.)}

Trưởng thành là đủ nội lực để không làm những điều mình sẽ hối hận.
\textit{(It is the inner strength not to do what you will regret.)}

\end{truyen}

\begin{baihoc}
Ranh giới quan trọng nhất không phải ranh giới với người khác, mà là ranh giới với phần yếu đuối và bốc đồng trong chính mình.
\textit{(The most important boundary is not with others, but with the impulsive and weaker part inside yourself.)}
\end{baihoc}

\begin{ghinhoanh}
\textbf{Short English to remember:}

The most important boundary is not with others, but with the impulsive and weaker part inside yourself.

\end{ghinhoanh}

\begin{tuhocnhanh}
\textbf{Cau hoi tu hoc / Self-study questions}

1. Van de chinh la gi? \textit{What is the main problem?}

2. Cach xu ly tot hon la gi? \textit{What is a better action?}

3. Mot cau tieng Anh em muon nho la cau nao? \textit{Which English sentence do you want to remember?}
\end{tuhocnhanh}
\ngancach

\section{Một Chữ “Không” Cứu Cả Đời}
\begin{truyen}{Một Chữ “Không” Cứu Cả Đời}{One No Can Save a Life}
\chuhoa{C}{ó những tai họa lớn bắt đầu từ việc một học sinh không dám nói không với chuyện rất nhỏ.}
\textit{(Some major disasters begin with a student’s inability to say no to something seemingly small.)}

Một lần chép bài hộ, một lần đi theo nhóm làm chuyện dại, một lần lên xe với người lái ẩu.
\textit{(One act of copying homework for someone, one outing to do something reckless, one ride with an unsafe driver.)}

Rất nhiều tai nạn bắt đầu kiểu đó.
\textit{(Many accidents begin that way.)}

Bản lĩnh không chỉ nằm ở những tình huống to tát.
\textit{(Strength is not found only in dramatic situations.)}

Nó nằm ở khoảnh khắc em ngắt được dòng cuốn của đám đông.
\textit{(It lives in the moment when you interrupt the momentum of the crowd.)}

Một chữ không nói ra đúng lúc nhiều khi đắt hơn một bài học nói ra sau khi đã hỏng.
\textit{(A no not spoken at the right time is often costlier than a lesson learned after damage has been done.)}

\end{truyen}

\begin{baihoc}
Dạy học sinh nói không không phải dạy cứng đầu. Đó là dạy năng lực tự bảo vệ trước sức ép xã hội.
\textit{(Teaching students to say no is not teaching stubbornness. It is teaching self-protection against social pressure.)}
\end{baihoc}

\begin{ghinhoanh}
\textbf{Short English to remember:}

Teaching students to say no is not teaching stubbornness.

It is teaching self-protection against social pressure.

\end{ghinhoanh}

\begin{tuhocnhanh}
\textbf{Cau hoi tu hoc / Self-study questions}

1. Van de chinh la gi? \textit{What is the main problem?}

2. Cach xu ly tot hon la gi? \textit{What is a better action?}

3. Mot cau tieng Anh em muon nho la cau nao? \textit{Which English sentence do you want to remember?}
\end{tuhocnhanh}
\ngancach

\section{Người Thợ Giữ Giới Hạn Của Mình}
\begin{truyen}{Người Thợ Giữ Giới Hạn Của Mình}{The Craftsman Who Kept His Limits}
\chuhoa{M}{ột người thợ nổi tiếng vì làm việc chắc tay, nhưng ông luôn từ chối nhận quá số việc mình có thể làm tử tế.}
\textit{(A craftsman became known for reliable work, yet always refused to take on more jobs than he could do well.)}

Người ta chê ông bỏ tiền.
\textit{(People said he was turning money away.)}

Ông đáp: ``Nhận quá sức là hại khách, hại nghề, và hại chính mình.''
\textit{(He replied, ``Taking on more than I can handle harms the client, the craft, and myself.'')}

Học sinh cũng vậy.
\textit{(Students are the same.)}

Gật đầu với mọi yêu cầu không làm em rộng lượng hơn nếu cuối cùng em làm đâu cũng dở và người cũng nát.
\textit{(Saying yes to every demand does not make you more generous if in the end everything is done poorly and you yourself are wrecked.)}

Ranh giới là thứ giữ chất lượng sống, không phải chỉ để làm khó người khác.
\textit{(Boundaries protect the quality of life; they are not merely obstacles placed before others.)}

\end{truyen}

\begin{baihoc}
Biết giới hạn của mình là trí khôn, không phải ích kỷ. Người không giữ giới hạn thường trả giá bằng sức khỏe và uy tín.
\textit{(Knowing your limits is wisdom, not selfishness. Those who fail to keep boundaries often pay with health and credibility.)}
\end{baihoc}

\begin{ghinhoanh}
\textbf{Short English to remember:}

Knowing your limits is wisdom, not selfishness.

Those who fail to keep boundaries often pay with health and credibility.

\end{ghinhoanh}

\begin{tuhocnhanh}
\textbf{Cau hoi tu hoc / Self-study questions}

1. Van de chinh la gi? \textit{What is the main problem?}

2. Cach xu ly tot hon la gi? \textit{What is a better action?}

3. Mot cau tieng Anh em muon nho la cau nao? \textit{Which English sentence do you want to remember?}
\end{tuhocnhanh}
\ngancach

\section{Câu Từ Chối Không Cần Gắt}
\begin{truyen}{Câu Từ Chối Không Cần Gắt}{A Refusal Does Not Need to Be Harsh}
\chuhoa{N}{hiều học sinh ngại từ chối vì nghĩ đã nói không là phải căng thẳng và làm mất lòng.}
\textit{(Many students avoid refusing because they think saying no must create tension and offense.)}

Thực ra có những câu rất đơn giản: ``Tao không làm chuyện đó,'' ``Tao không đi đâu,'' ``Tao đang quá tải rồi.''
\textit{(In fact, there are simple sentences: ``I am not doing that,'' ``I am not going,'' ``I am already overloaded.'')}

Ranh giới rõ mà bình tĩnh thường có sức nặng hơn kiểu giải thích dài dòng xin phép được tồn tại.
\textit{(Clear and calm boundaries often carry more weight than lengthy justifications begging permission to exist.)}

Học sinh cần luyện câu từ chối như luyện bài tập.
\textit{(Students need to practice refusal as they practice exercises.)}

Đến lúc áp lực đến, thứ đã luyện mới bật ra được.
\textit{(When pressure comes, only what has been practiced will come out.)}

\end{truyen}

\begin{baihoc}
Dám nói không không đồng nghĩa với dữ dằn. Nhiều ranh giới vững nhất được nói bằng giọng rất bình thường.
\textit{(Having the courage to say no does not mean being aggressive. Many of the strongest boundaries are spoken in an ordinary tone.)}
\end{baihoc}

\begin{ghinhoanh}
\textbf{Short English to remember:}

Having the courage to say no does not mean being aggressive.

Many of the strongest boundaries are spoken in an ordinary tone.

\end{ghinhoanh}

\begin{tuhocnhanh}
\textbf{Cau hoi tu hoc / Self-study questions}

1. Van de chinh la gi? \textit{What is the main problem?}

2. Cach xu ly tot hon la gi? \textit{What is a better action?}

3. Mot cau tieng Anh em muon nho la cau nao? \textit{Which English sentence do you want to remember?}
\end{tuhocnhanh}
\ngancach

\section{Không Thích Không Có Nghĩa Phải Cố Hợp}
\begin{truyen}{Không Thích Không Có Nghĩa Phải Cố Hợp}{Disliking Something Does Not Mean You Must Force Yourself to Fit}
\chuhoa{C}{ó học sinh nhận mọi lời mời, mọi vai trò, mọi hoạt động chỉ vì sợ mình khó gần.}
\textit{(Some students accept every invitation, role, and activity because they fear seeming difficult.)}

Nhưng con người không lớn lên bằng cách luôn ép mình hợp với mọi chỗ.
\textit{(But people do not mature by constantly forcing themselves to fit everywhere.)}

Có việc em không thích đơn giản vì nó không hợp em.
\textit{(Some things you dislike simply do not suit you.)}

Không cần biến điều đó thành lỗi đạo đức.
\textit{(There is no need to turn that fact into a moral defect.)}

Biết nói không với cái không hợp cũng là cách giữ năng lượng cho cái đáng.
\textit{(Knowing how to refuse what does not fit is also a way of preserving energy for what matters.)}

\end{truyen}

\begin{baihoc}
Học sinh không cần vừa lòng tất cả để trở thành người tử tế. Biết chọn cũng là một phần của trưởng thành.
\textit{(Students do not need to please everyone in order to be decent. Selection is part of maturity too.)}
\end{baihoc}

\begin{ghinhoanh}
\textbf{Short English to remember:}

Students do not need to please everyone in order to be decent.

Selection is part of maturity too.

\end{ghinhoanh}

\begin{tuhocnhanh}
\textbf{Cau hoi tu hoc / Self-study questions}

1. Van de chinh la gi? \textit{What is the main problem?}

2. Cach xu ly tot hon la gi? \textit{What is a better action?}

3. Mot cau tieng Anh em muon nho la cau nao? \textit{Which English sentence do you want to remember?}
\end{tuhocnhanh}
\ngancach

\section{Từ Chối Một Người Lớn}
\begin{truyen}{Từ Chối Một Người Lớn}{Refusing an Adult}
\chuhoa{K}{hó nhất không phải lúc từ chối bạn bè.}
\textit{(The hardest refusals are not always to peers.)}

Khó nhất là khi người ép em là một người lớn có quyền hơn.
\textit{(Sometimes the hardest moment comes when the pressure is from an adult with more power.)}

Một học sinh bị bắt làm thêm việc tập thể vô lý trong lúc đã quá tải học hành.
\textit{(A student was forced into unreasonable extra duties while already overloaded with study.)}

Em rất sợ nói không vì sợ bị ghét.
\textit{(He feared saying no because he feared being disliked.)}

Cuối cùng em học cách nói: ``Em muốn giúp, nhưng hiện tại em không đảm bảo làm tốt được.''
\textit{(Eventually he learned to say, ``I want to help, but at present I cannot guarantee I will do it well.'')}

Từ chối người có quyền cần nhiều khôn ngoan hơn, nhưng vẫn là kỹ năng cần có.
\textit{(Refusing someone with authority requires more tact, but it is still a necessary skill.)}

\end{truyen}

\begin{baihoc}
Học sinh cần được dạy cách bảo vệ ranh giới cả trước người trên, không chỉ với bạn bè ngang hàng.
\textit{(Students need to learn how to protect boundaries not only with peers, but also with authority figures.)}
\end{baihoc}

\begin{ghinhoanh}
\textbf{Short English to remember:}

Students need to learn how to protect boundaries not only with peers, but also with authority figures.

\end{ghinhoanh}

\begin{tuhocnhanh}
\textbf{Cau hoi tu hoc / Self-study questions}

1. Van de chinh la gi? \textit{What is the main problem?}

2. Cach xu ly tot hon la gi? \textit{What is a better action?}

3. Mot cau tieng Anh em muon nho la cau nao? \textit{Which English sentence do you want to remember?}
\end{tuhocnhanh}
\ngancach

\section{Khi Em Nói Không Với Phiên Bản Tệ Của Mình}
\begin{truyen}{Khi Em Nói Không Với Phiên Bản Tệ Của Mình}{When You Say No to the Worse Version of Yourself}
\chuhoa{C}{ó lúc ranh giới lớn nhất không nằm bên ngoài.}
\textit{(At times the biggest boundary is not outside.)}

Nó nằm trong đầu khi em muốn buông xuôi, muốn trả đũa, muốn nói dối cho xong.
\textit{(It lies in the mind when you want to give up, retaliate, or lie to escape.)}

Không ai nhìn thấy những lần nói không với bản thân đó, nhưng nó tạo nên phần cứng nhất của nhân cách.
\textit{(No one sees those refusals inside the self, yet they build the hardest part of character.)}

Người giữ được mình trong lúc không ai nhìn thường là người đứng được lâu nhất khi ai cũng nhìn.
\textit{(The person who can govern themselves when no one watches is often the one who stands longest when everyone watches.)}

Tự chủ không hào nhoáng, nhưng nó là xương sống của tự do.
\textit{(Self-control is not glamorous, but it is the backbone of freedom.)}

\end{truyen}

\begin{baihoc}
Cuối cùng, chữ không quan trọng nhất là chữ em nói với phần thấp của chính mình. Từ đó mọi ranh giới khác mới có gốc.
\textit{(In the end, the most important no is the one you say to the lower part of yourself. From there all other boundaries gain roots.)}
\end{baihoc}

\begin{ghinhoanh}
\textbf{Short English to remember:}

In the end, the most important no is the one you say to the lower part of yourself.

From there all other boundaries gain roots.

\end{ghinhoanh}

\begin{tuhocnhanh}
\textbf{Cau hoi tu hoc / Self-study questions}

1. Van de chinh la gi? \textit{What is the main problem?}

2. Cach xu ly tot hon la gi? \textit{What is a better action?}

3. Mot cau tieng Anh em muon nho la cau nao? \textit{Which English sentence do you want to remember?}
\end{tuhocnhanh}
